\section{Native contact models: configurations and results of our configurations}\label{sup:natives}

% The paper cites these subsections by number (S3.1 to S3.6): keep their order.

Everything in this section is a result of the configurations described in
Section~\ref{sup:native-config}. They are ours: none was reviewed by the developers of an
engine. The paper draws from the native runs only what follows from a property of a model
family; the results below go beyond that, several are unexplained, and none should be read as
the performance of an engine in expert hands.

\subsection{Configurations}\label{sup:native-config}

Each native model has one configuration, the same for every scenario, built from the engine's
documentation and written down with its sources in the specification of the repository
(\code{docs/scenarios.md}, Section~3). The names are used as written here throughout;
``point-contact models'' means the five other than Drake hydroelastic, which are a penalty
model with one point per pair (Drake point), a penalty model on a manifold (Chrono SMC), a
complementarity solver (Chrono NSC) and soft constraints (the two of MuJoCo). Every native run
has a wall-clock limit of one hour in the base variant and ten minutes in the series (half an
hour for the methods of the library).

The scenario gives every method the same $K$, $h$ and $\mu$. Models whose stiffness belongs to
a contact point (N/m) have no areal stiffness on which to map $K$; they use a reference area
$A_\mathrm{ref}$, one sixth of the surface of a solid (the face of a cube), so that
$K A_\mathrm{ref}$ is the stiffness of a foundation $K$ under a face of that size. This sets
the scale and is exact only when the contact is a face of that area.

\paragraph{Drake, hydroelastic} (the default native model of Drake in this paper)
\cite{DrakeHydroGuide,DrakeHydroParams}. Planar grounds are compliant half-spaces and other
grounds are compliant tetrahedral meshes, with hydroelastic modulus $E = K H$, $H$ being the
depth of the pressure field, so that the pressure is $K$ times the depth: with a rigid
counterpart this is a Winkler foundation of the same $K$ evaluated on the rigid surface. Bodies
are rigid unless they touch another rigid body; the contact graph is two-coloured from the
grounds. A rigid body that touches a half-space has its surface refined until no edge exceeds
one eighth of its largest side, because Drake turns each polygon of the contact surface into one
contact point and the guide warns that the refinement of a contact with a half-space is that of
the other geometry; an unrefined box stands on two points. Dissipation is of Hunt--Crossley
type with coefficient $h$, with the discrete approximation that keeps that form (``lagged''),
and the friction coefficients of the geometries are solved so that Drake's combination rule
gives the $\mu$ of each pair.

\paragraph{Drake, point} \cite{DrakePointContact}. Compliant point contact in the same plant
and solver, between analytic primitives or convex pieces, with point stiffness $K A_\mathrm{ref}$
per geometry (Drake combines two geometries in series) and dissipation $h$. A box carries a
small sphere at each corner, the device of Drake's own box models \cite{DrakeFoamBrick}, since
one point per pair cannot hold a face.

\paragraph{Chrono, SMC} (the default for Chrono) \cite{ChronoCollisions,ChronoSMCSource}.
Analytic primitives where the scenario has them and convex pieces otherwise. Hooke's law with
explicit coefficients: normal and tangential stiffness $K A_\mathrm{ref}/4$ per contact point
(a face contact has a manifold of four), and a normal damping that is the Hunt--Crossley term
linearized at the static penetration under the body's weight. In the version used, the
Jacobians of the contact force are disabled in the source and the tangential displacement is
that of one step, so the model is explicit and has no stiction memory; the backend sub-steps so
that the contact oscillation is resolved with at least 60 steps, and uses the stiffest tangential
damping that an explicit step integrates, which bounds the creep of a held body but does not
remove it.

\paragraph{Chrono, NSC} \cite{ChronoSystem}. The complementarity formulation for rigid contact,
with the same shapes, the pair's $\mu$ and zero restitution; $K$ and $h$ do not enter. Solved
with the Barzilai--Borwein method \cite{Barzilai1988}, which the manual recommends when
accuracy matters.

\paragraph{MuJoCo, soft} (the default for MuJoCo) \cite{MuJoCoDocs,MuJoCoModeling}. MuJoCo's
contact model with the values that its manual recommends for new models: analytic primitives
or convex pieces, elliptic friction cones, a tight solver tolerance, a few no-slip iterations
and the \code{implicitfast} integrator. The stiffness of the constraint is set so that the
static force of a flat contact is $K A_\mathrm{ref}$ times the penetration; only the statics
are a spring's, because the constraint also removes a fixed share of the load that presses the
bodies together. Damping is critical: MuJoCo's damper is linear and cannot follow the dependence
of the Hunt--Crossley restitution on the impact speed. Each contact has torsional friction with
the coefficient of a disk of its patch area under uniform pressure.

\paragraph{MuJoCo, default.} The same geometry with MuJoCo's own default parameters and
options: what a user obtains without changing anything. $K$ and $h$ do not enter.

\subsection{Revisions and known asymmetries}

The configurations were not fixed blind: after the first complete run of the development
scenarios several options were revised with the errors of the battery in view. The revisions
were the refinement of rigid surfaces to one eighth of the largest side for Drake's
hydroelastic contact (\scn{block} went from 13 to 1.4\,\%), the corner spheres of boxes for
Drake's point contact (\scn{block} from 98 to 0.1\,\%), the solver tolerance and the reference
area for MuJoCo (\scn{stack} from 33 to 0.16\,\%) and the choice of the Barzilai--Borwein solver
for Chrono NSC. Each revision was a setting applied to every scenario and was chosen from a
sweep over all of them, never per scenario; they are nonetheless iterations on the scores of
the development set, fewer than those of the laws, recorded in the history of the repository
(\code{docs/history.md}, entries of 11 November 2026). % The four figures are those of that history, single runs on the development scenarios.

Drake's point contact gives a box a small sphere at each corner and
gives a cylinder or a disk nothing: its flat face touches the ground at one point, so the errors
and the runs without a stop time of that model on \scn{disk}, \scn{contensou} and \scn{cone} are
in part a consequence of our configuration and not of the model. Chrono SMC is sub-stepped by the
backend (between 1 and 15 sub-steps per step of the battery), so its step is not the step of the
other methods, in the convergence series and in the costs. The reference area $A_\mathrm{ref}$ is
a scale: for bodies that are not cubes (the plates of \scn{offset} and \scn{tower}, the disk) the
point stiffness of a native model differs from that of the foundation under the face in contact by
factors of 2 to 3. The static penetration of each native model was not compared with the
foundation's scenario by scenario. MuJoCo's recommended configuration misses \scn{block}, the
simplest scenario, by \result{err/block/soft} (\result{pertrange/block/soft/range} over the
perturbed starts); a configuration reviewed by a user of that engine would probably not, and the
result should not be read as representative of expert use.

\subsection{Results by class}

The errors are in the tables of the paper (Tables~3 to~5) and in Section~\ref{sup:errors}.

\emph{Class A.} On the seven scenarios of a body sliding or rolling on one plane, Drake point
and Chrono NSC are as accurate as the two laws or more: in
\result{within/oneplane/dpoint+nsc/0.4} of them the errors of both models are below 0.4\,\%,
though Drake point misses \scn{vgroove} by \result{err/vgroove/dpoint}. MuJoCo soft has
\result{err/block/soft}, \result{err/twoblocks/soft} and \result{err/billiard/soft} in
\scn{block}, \scn{twoblocks} and \scn{billiard}, and MuJoCo default \result{err/push/mjdefault}
and \result{err/twoblocks/mjdefault} in \scn{push} and \scn{twoblocks}. Chrono SMC, which has
no stick memory in the version used, creeps in \scn{push} and \scn{wedge}
(\result{err/push/smc}, \result{err/wedge/smc}). In \scn{bowl} and \scn{trough} the native
models lose between \result{err/bowl/smc} and \result{err/trough/soft} of the amplitude. In
\scn{turntable} Drake point keeps the orbit, with \result{err/turntable/dpoint} and
\result{pertrange/turntable/dpoint/range} over the perturbations; the base run of Chrono NSC
(\result{err/turntable/nsc}) is the best of its perturbed starts
(\result{pertrange/turntable/nsc/range}, median \result{pertrange/turntable/nsc/median}). In
\scn{journal}, Drake hydroelastic and MuJoCo soft are within 5\,\% (\result{err/journal/hydro},
\result{err/journal/soft}; \result{pertrange/journal/soft/range} when perturbed), and Drake
point and Chrono's two models are between \result{err/journal/smc} and
\result{err/journal/nsc}. On the cut ground Chrono's models depend on the cut
(\result{err/pieces64/smc}, and \result{err/pieces2/nsc} for NSC already on two pieces). In
\scn{wedge}, Chrono NSC, MuJoCo soft and Drake point hold the block that must hold. In
\scn{oblique} Chrono SMC and Drake point, on a sphere, are within 0.6\,\%. Chrono NSC and MuJoCo
soft hold \scn{tower} and \scn{massratio} within 1\,\%; Drake hydroelastic holds \scn{massratio}
and did not finish \scn{tower}; Drake point, which uses the same solver, loses both
(\result{err/massratio/dpoint} and \result{err/tower/dpoint}), so being implicit in the contact
force is not sufficient. In \scn{energy} every native model but Chrono SMC has an error above
20\,\% (Drake point \result{err/energy/dpoint}).

\emph{Class B.} Drake point, Chrono NSC and the two MuJoCo models are each outside the band in
\result{band/dpoint/outside} of the five scenarios; Chrono SMC is outside in
\result{band/smc/outside} and has no stop time in \result{band/smc/none}
(Table~\ref{tab:band}). Their errors against the nominal value go from \result{err/ring/smc}
(Chrono SMC in \scn{ring}, below the \result{err/cone/hydro} of Drake hydroelastic in
\scn{cone}) to \result{err/bearing/nsc}. Drake hydroelastic is inside the band in
\result{band/hydro/inside} scenarios and undetermined in the others, with \result{err/disk/hydro}
in \scn{disk}.

\emph{Class C.} Drake hydroelastic is within 5\,\% in \result{pass/C/hydro/5}; \scn{tip} ends
without a value with every point-contact model, as designed: a box on stiff points tips at the
rigid angle, which the ramp of the scenario never reaches.

\emph{Step, mesh and stiffness.} Drake hydroelastic is the method whose error depends most on
the step: it converges in \result{convn/dt/hydro/converging} series, and its errors at 1\,ms in
the rolling and impact scenarios are step errors. It holds \scn{massratio} at every step at
which it finished, as do Chrono NSC and the two MuJoCo models. Chrono SMC is sub-stepped by the
backend, so its series is not one of the step of the battery. With $K\times10$ the errors of
Drake hydroelastic on class A fall in \result{stiffn/hydro/better} of \result{stiffn/hydro/n}
scenarios, and on the flat \scn{disk} its error does not rise (\result{err/disk/hydro/stiff}).

\subsection{Scatter under perturbation}\label{sup:scatter}

Over the eight perturbations of the initial velocities (1\,mm/s and 0.01\,rad/s, from a seeded
generator), the range of the error relative to its median has a median over the scenarios of
\result{pert/hydro} for Drake hydroelastic and \result{pert/soft} for MuJoCo soft, and of
\result{pert/dpoint}, \result{pert/nsc} and \result{pert/smc} for Drake point, Chrono NSC and
Chrono SMC (for the two laws, \result{pert/gonthier} and \result{pert/distributed}). The base
errors are therefore representative of nearby initial conditions for Drake hydroelastic and for
MuJoCo soft, and less so for Drake point and the two Chrono models.
Table~\ref{tab:perturbations} lists the cells of classes A and B in which the perturbed runs of
a native model differ by more than a factor of 2 or lose a run: for those, the base value is
one draw and the median is the better summary. Chrono NSC in \scn{disk} has
\result{err/disk/nsc} in the base run and \result{pertrange/disk/nsc/range} over the
perturbations (median \result{pertrange/disk/nsc/median}).

\begin{table}[!htbp]
\caption{Native models whose error scatters over the eight perturbed starts (classes A and B):
error [\%] of the base run, median and range of the perturbed runs that have a value, and number
of perturbed runs without one}
\label{tab:perturbations}
\generatedtable{perturbations}
\end{table}

\subsection{Unexplained results}

The following results of our configurations have no explanation, and any of them may be an
effect of the configuration. Drake hydroelastic: \scn{hole} (\result{err/hole/hydro}),
\scn{corner} (\result{err/corner/hydro}) and \scn{stack} (\result{err/stack/hydro}), none of
which falls with the step; \scn{tower} and the smallest steps of \scn{stack} and
\scn{massratio} reach the time limit. Chrono NSC: \scn{oblique} and most of \scn{restitution}
end without a rebound value (the model was given zero restitution, so it was not expected to
meet the rebounds, but why some are not produced at all was not traced), \scn{edge} has
\result{err/edge/nsc}, and its mesh variants on \scn{disk} and \scn{turntable}
(\result{err/disk/nsc/mesh32} and \result{err/turntable/nsc/mesh32} at 32 segments,
\result{err/disk/nsc/mesh128} and \result{err/turntable/nsc/mesh128} at 128) differ from its
base runs (\result{err/disk/nsc}, \result{err/turntable/nsc}) although its shapes are analytic
primitives; the scatter of the model over nearby starts is of that size, and what else differs
between those runs was not traced. MuJoCo default: \scn{edge} (\result{err/edge/mjdefault}).
The runs of \scn{disk} and \scn{contensou} with Chrono SMC and Drake point, and of \scn{cone}
with the latter, end without a stop time: in the diagnostic runs of Chrono SMC the spinning
body wobbled on its rim and kept a spin above the 1\,\% that defines the stop; Drake's runs
were not traced, and its disk stands on one point by our configuration.

\subsection{A choice of model by situation, as configured}

Table~\ref{tab:decision-natives} adds to the table of the paper's discussion what the runs of
our configurations show about particular native models.

\begin{table}[!htbp]
\caption{Native models by situation, as far as our configurations show (classes A and B, base
variant)}
\label{tab:decision-natives}
\footnotesize
\begin{tabular}{@{}p{0.27\linewidth}p{0.27\linewidth}p{0.40\linewidth}@{}}
\toprule
Situation & Native models that did well & Evidence \\
\midrule
A rigid body sliding or rolling on one plane; sphere and cylinder contacts &
  Drake point, Chrono NSC &
  \scn{block}, \scn{sphere}, \scn{cylinder}, \scn{twoblocks}: Drake point
  \result{err/block/dpoint}, \result{err/sphere/dpoint}, \result{err/cylinder/dpoint},
  \result{err/twoblocks/dpoint}; \law{distributed} \result{err/block/distributed},
  \result{err/sphere/distributed}, \result{err/cylinder/distributed},
  \result{err/twoblocks/distributed} \\
A flat face that spins, or slides and spins &
  Drake hydroelastic &
  \scn{disk}: \result{err/disk/hydro}; \law{distributed} \result{err/disk/distributed};
  point-contact models outside the band or scattered (Tables~\ref{tab:band}
  and~\ref{tab:perturbations}) \\
Curved conforming contact &
  Drake hydroelastic &
  \scn{bearing} at 128 segments: \result{err/bearing/hydro/mesh128}; \law{pcm}
  \result{err/bearing/pcm/mesh128}, \law{distributed} \result{err/bearing/distributed/mesh128} \\
A body held by friction near its limit &
  Chrono NSC, MuJoCo soft, Drake point &
  \scn{wedge}: \result{err/wedge/nsc}, \result{err/wedge/soft}, \result{err/wedge/dpoint};
  Chrono SMC \result{err/wedge/smc} \\
Stiff contacts on light bodies, large mass ratios, tall stacks at large steps &
  Chrono NSC, MuJoCo soft; Drake hydroelastic where its cost allows. Not every implicit solver &
  \scn{massratio} at 1\,ms: Chrono NSC, MuJoCo soft and Drake hydroelastic within
  \result{err/massratio/soft}; Drake point \result{err/massratio/dpoint}, and
  \result{err/tower/dpoint} in \scn{tower} \\
Long rolling &
  Drake point &
  \scn{turntable}: \result{err/turntable/dpoint} (\result{pertrange/turntable/dpoint/range}
  perturbed); the laws \result{err/turntable/distributed} \\
\botrule
\end{tabular}
\end{table}
